\documentclass[10pt,a4paper]{article}

%double si on veut une version prof et une version élève. simple si on veut une seule version.
\newcommand{\buildmode}{simple}

\InputIfFileExists{version.tex}{}{}
% Version (définie par le script)
\providecommand{\version}{eleve}

\usepackage{feuilleinterro}
\usepackage{schemaevolution}

%%%à modifier à chaque fois

\renewcommand{\niveau}{1M}
\renewcommand{\chapitre}{Évolutions}
\renewcommand{\numchap}{0.3}
\renewcommand{\theme}{Statistiques}

\begin{document}

% =========================================================
% PAGE 1 : VERSION A
% =========================================================

\interroversion{A}

\textbf{Exercice 1 -- Solution d'une équation \hfill /2}

Le nombre \(3\) est-il solution de l'équation \(2x+3=5x-6\) ? Justifier.

\vspace{4cm}

\textbf{Exercice 2 -- Résolution d'une équation \hfill /3}

Résoudre l'équation \(6x+5=9\). Donner la solution sous la forme d'une fraction irréductible.

\vspace{4cm}

\textbf{Exercice 3 -- Valeur finale \hfill /2}

Compléter le schéma suivant, en justifiant par un calcul.

\schemaevolution{60}{\dots}{-20\%}{}

\vspace{3.5cm}

\textbf{Exercice 4 -- Taux d'évolution \hfill /3}

Compléter le schéma suivant, en justifiant par un calcul.

\schemaevolution{80}{100}{\dots}{}

\newpage

% =========================================================
% PAGE 2 : VERSION B
% =========================================================

\interroversion{B}

\textbf{Exercice 1 -- Solution d'une équation \hfill /2}

Le nombre \(4\) est-il solution de l'équation \(3x-2=x+5\) ? Justifier.

\vspace{4cm}

\textbf{Exercice 2 -- Résolution d'une équation \hfill /3}

Résoudre l'équation \(4x-3=7\). Donner la solution sous la forme d'une fraction irréductible.

\vspace{4cm}

\textbf{Exercice 3 -- Valeur finale \hfill /2}

Compléter le schéma suivant, en justifiant par un calcul.

\schemaevolution{40}{\dots}{+25\%}{}

\vspace{3.5cm}

\textbf{Exercice 4 -- Taux d'évolution \hfill /3}

Compléter le schéma suivant, en justifiant par un calcul.

\schemaevolution{50}{40}{\dots}{}

\newpage

% =========================================================
% PAGE 3 : VERSION C
% =========================================================

\interroversion{C}

\textbf{Exercice 1 -- Solution d'une équation \hfill /2}

Le nombre \(5\) est-il solution de l'équation \(4x-7=2x+3\) ? Justifier.

\vspace{4cm}

\textbf{Exercice 2 -- Résolution d'une équation \hfill /3}

Résoudre l'équation \(10x+1=5\). Donner la solution sous la forme d'une fraction irréductible.

\vspace{4cm}

\textbf{Exercice 3 -- Valeur finale \hfill /2}

Compléter le schéma suivant, en justifiant par un calcul.

\schemaevolution{80}{\dots}{-25\%}{}

\vspace{3.5cm}

\textbf{Exercice 4 -- Taux d'évolution \hfill /3}

Compléter le schéma suivant, en justifiant par un calcul.

\schemaevolution{200}{250}{\dots}{}

\newpage

% =========================================================
% PAGE 4 : VERSION D
% =========================================================

\interroversion{D}

\textbf{Exercice 1 -- Solution d'une équation \hfill /2}

Le nombre \(2\) est-il solution de l'équation \(5x+1=3x+4\) ? Justifier.

\vspace{4cm}

\textbf{Exercice 2 -- Résolution d'une équation \hfill /3}

Résoudre l'équation \(8x-5=1\). Donner la solution sous la forme d'une fraction irréductible.

\vspace{4cm}

\textbf{Exercice 3 -- Valeur finale \hfill /2}

Compléter le schéma suivant, en justifiant par un calcul.

\schemaevolution{50}{\dots}{+20\%}{}

\vspace{3.5cm}

\textbf{Exercice 4 -- Taux d'évolution \hfill /3}

Compléter le schéma suivant, en justifiant par un calcul.

\schemaevolution{20}{15}{\dots}{}

\end{document}
